The correlation values in the matrices represent Pearson correlation coefficients calculated between pairs of models based on their attack success rates across the 30 test cases in batch7. For each model pair, we compute the correlation between their respective score vectors across all test cases, where single-turn scenarios use the single score and multi-turn scenarios use the maximum score achieved across all turns. A correlation of 1.0 indicates that two models respond identically to the same attacks (perfect positive correlation), while a correlation of -1.0 would indicate completely opposite responses (perfect negative correlation). A correlation near 0 suggests that the models' vulnerabilities are largely independent.

The analysis reveals interesting patterns in model behavior: models from the same laboratory tend to exhibit higher correlations (highlighted with thicker borders), suggesting shared training methodologies or safety measures. Multi-turn scenarios generally show slightly higher correlations than single-turn scenarios, indicating that models may converge toward similar response patterns when given multiple interaction rounds. The correlation matrices provide insights into which models might be redundant for evaluation purposes and which offer distinct vulnerability profiles for comprehensive safety testing.